\section{A uniform estimate at one volume}\label{sec:apex}

The preceding surface estimates control cohomology in terms of a
discriminant, without a factor depending on the rank.  We now use them to
prove the following estimate at one fixed value of the volume parameter.
The factor on its right measures the distance from the equality case of
the first-tilt discriminant inequality.

\begin{proposition}\label{prop:apex}
Let $X$ be a smooth projective complex threefold, let $H$ be very ample,
and let $B_0$ be a rational divisor class with $H^2B_0=0$.
There are constants $A>0$ and $C\geq0$, depending only on $X,H,B_0$,
such that, for every $b\in\RR$ and every finite-slope
$\nu_{A,b}$-semistable object $E\in\mathcal B_{H,B_0+bH}$
with $\nu_{A,b}(E)=0$,
\begin{equation}\label{eq:apex-estimate}
 z_B(E)-\frac{A^2}{6}d_B(E)
 \leq C\bigl(d_B(E)-A|r(E)|\bigr),\qquad B=B_0+bH.
\end{equation}
The same constants can be chosen for $B_0$ and $-B_0$.
\end{proposition}

Here $h=H^3$ and $d_B,w_B,z_B$ are normalized by $h$ as in
Section~\ref{sec:tilt}.  In particular, $d_B\geq A|r|$ at slope zero.
We first choose all the constants from the fixed geometry.  For an
individual $E$, we then separate a narrow range of sheaf slopes from a
remainder whose cohomology is controlled by
\[
 e=d_B(E)-A|r(E)|.
\]
Reduction to characteristic $p$ and Riemann--Roch convert these
cohomology estimates into \eqref{eq:apex-estimate}.  The remainder
contributes at most $Cp^3e$.  The narrow sheaf tail can have large
rank even when $e=0$; two restriction sequences will instead bound
its sections by a cubic expression in its slopes, producing the
coefficient $A^2/6$.  Only errors that vanish after division by
$p^3$ will be allowed to depend on $E$.
We write $F(t)=F\otimes\OO_X(tH)$, also for complexes, and write
$h^i(F)=\dim\HH^i(X,F)$.

\subsection{Constants chosen before the object}

Tensoring by $\OO_X(kH)$ identifies the parameters $b$ and $b+k$
and preserves all the $B$-twisted characters.  It therefore suffices
to treat $|b|\leq1$.  Choose a positive integer $q$ and a divisor $D$
representing $qB_0$.  On a smooth reduction of $X$ in characteristic
$p$, let $F$ denote absolute Frobenius and put
\begin{equation}\label{eq:rounding-lines}
 L_p=\OO_X\bigl(-\lfloor p/q\rceil D-\lfloor pb\rceil H\bigr),
 \qquad L'_p=K_X\otimes L_p^{-1}.
\end{equation}
Either nearest integer can be chosen at a tie.  Thus
$c_1(L_p)/p\longrightarrow-B$.
We will estimate Frobenius sections through sections at two fixed
twists on the original object.  The following presentation supplies
those twists; its consequence immediately below explains their roles.

\begin{lemma}\label{lem:frobenius-presentation}
There are positive integers $j,j_2$ and a constant $C_0$, determined
by $X,H,D$, with the following property.  After choosing and shrinking
a model of this fixed geometry over a finitely generated integral
subring of $\CC$, every smooth geometric fiber in characteristic $p$
and every $|b|\leq1$ satisfy the following assertion, for either
$L=L_p$ or $L=L'_p$:
\[
 G_{p,L}=F_*(K_X^{1-p}\otimes L^{-1})
\]
has a presentation
\begin{equation}\label{eq:frobenius-presentation}
 0\longrightarrow K_{p,L}\longrightarrow\OO_X(-j)^{N_{p,L}}
 \longrightarrow G_{p,L}\longrightarrow0,
 \qquad N_{p,L}\leq C_0p^3,
\end{equation}
where $K_{p,L}(j_2)$ is generated by its global sections.
\end{lemma}

\begin{proof}
For each line bundle in the fixed list
$\OO_X(\pm D),\OO_X(\pm H),K_X^{\pm1}$, choose three successive
surjections from sums of negative powers of $\OO_X(H)$, with locally
free kernels.  They exist by global generation after twisting; a
surjection onto a locally free sheaf has locally free kernel.
Spread these finitely many sequences and a regularity bound for
$\OO_X$ to a model, and shrink so that they remain exact on the
fibers.

There is consequently an integer $v_0$, independent of the fiber,
such that tensoring by any one line bundle in the list increases by
at most $v_0$ a threshold above which all higher cohomology vanishes.
Indeed, tensor its three-step presentation by the original line
bundle and use the cohomology sequences.  The three dimension shifts
end in cohomological degree greater than three.  If the finitely many
twists in the presentation are bounded by $v_0$, every intervening
sum of lines has vanishing higher cohomology above the shifted
threshold.  Starting with the fixed bound for $\OO_X$ proves the
claim by induction on the number of tensor factors.

The line $K_X^{1-p}\otimes L^{-1}$ is a product of at most $c_0p+c_1$
members of this fixed list, uniformly for $|b|\leq1$ and for both
choices of $L$.  It therefore has no higher cohomology after twisting
by $sH$ whenever $s\geq c_2p+c_3$.  Projection formula gives
\[
 H^i\bigl(G_{p,L}(j-i)\bigr)
 =H^i\bigl(K_X^{1-p}\otimes L^{-1}(p(j-i))\bigr)=0
 \quad(1\leq i\leq3)
\]
for a fixed sufficiently large $j$.  Castelnuovo--Mumford regularity
then gives the evaluation surjection in
\eqref{eq:frobenius-presentation}, using a basis of
$H^0(G_{p,L}(j))$, and gives surjectivity of multiplication of these
sections at every higher twist.

The cohomology sequence for its kernel gives a uniform regularity
bound $j_2$.  For degrees at least two use regularity of $G_{p,L}$
and of $\OO_X(-j)$; in degree one use the just-proved multiplication
surjectivity.  Increasing $j_2$ makes the kernel generated at $j_2$.
Finally, the higher cohomology of
$K_X^{1-p}\otimes L^{-1}(pj)$ vanishes, so
\[
 N_{p,L}=\chi\bigl(K_X^{1-p}\otimes L^{-1}(pj)\bigr).
\]
This Euler characteristic is constant under specialization for
each of the indicated line bundles.  On the complex fiber,
Riemann--Roch is a polynomial of degree at most three in exponents
bounded by a fixed multiple of $p$.  This proves $N_{p,L}\leq C_0p^3$.
All the choices preceded any choice of a tilt object.
\end{proof}

\begin{corollary}[From fixed twists to Frobenius sections]
\label{cor:frobenius-section-transfer}
On a smooth geometric fiber in the preceding lemma, let $P_1$ be a
perfect complex with ordinary cohomology in degrees $[-1,0]$.
If
\[
 H^0\bigl(\mathcal H^{-1}(P_1)(j_2)\bigr)=0,
\]
then, for either $L=L_p$ or $L=L'_p$,
\begin{equation}\label{eq:frobenius-section-transfer}
 h^0(F^*P_1\otimes L)
 \leq C_0p^3h^0(P_1(j)).
\end{equation}
The twists $j,j_2$ and the constant $C_0$ are those chosen from the
fixed geometry in Lemma~\ref{lem:frobenius-presentation}.
\end{corollary}

\begin{proof}
Finite duality and projection formula give
\begin{equation}\label{eq:frobenius-hom-identity}
 h^0(F^*P_1\otimes L)=\dim\Hom(G_{p,L},P_1).
\end{equation}
Indeed $F^!P_1=F^*P_1\otimes K_X^{1-p}$, so adjunction applied
to $K_X^{1-p}\otimes L^{-1}$ proves the identity.
The semilinearity of absolute Frobenius does not change these
dimensions.  Apply \eqref{eq:frobenius-presentation}.  The obstruction
to injectivity of
\[
 \Hom(G_{p,L},P_1)\longrightarrow
 \Hom(\OO_X(-j)^{N_{p,L}},P_1)
\]
is the image of
\[
 \Hom(K_{p,L}[1],P_1)
 =\Hom(K_{p,L},\mathcal H^{-1}(P_1)).
\]
The equality uses the ordinary cohomological range of $P_1$.
Generation of $K_{p,L}(j_2)$ and the assumed vanishing make this
group zero.  The Hom map is therefore injective, and its target
has dimension $N_{p,L}h^0(P_1(j))\leq C_0p^3h^0(P_1(j))$.
\end{proof}

Thus the characteristic-zero estimates we shall need are a section
bound at $j$ and a vanishing for the negative cohomology sheaf at
$j_2$.  They will pass to the chosen fibers by semicontinuity.
We now choose the remaining constants so that both estimates can
be proved uniformly in the tilt object.

Fix $m$ and the compact smooth neighborhoods of surfaces and curves
from Lemma~\ref{lem:uniform-sections}.  In particular,
\begin{equation}\label{eq:apex-m-choice}
 h^0(\OO_S(m))=\chi(\OO_S(m)),\qquad
 \mu(K_S)-m\leq-1.
\end{equation}
Put $M=\max(j,j_2)$, choose $c>M+2$, and then choose $A$ so large that
Proposition~\ref{prop:negative-surface-cohomology} applies for upper
twist $M$ and both $B_0,-B_0$, and
\begin{equation}\label{eq:apex-A-choice}
 A>c,\qquad A>j_2+2,\qquad A\geq1,\qquad
 A^2-c^2\geq\delta_{\rm line}:=-HB_0^2/h.
\end{equation}
The number $\delta_{\rm line}$ is nonnegative by the Hodge index
theorem.  We shall choose one further integer $m_0$ below, depending
only on these constants, and put $t_0=-m_0$.  Thus $A$ does not
depend on the negative endpoint $t_0$: the negative-surface estimate
requires an upper bound on the twists for its positivity, and its
constant may subsequently depend on the chosen finite range.

\subsection{The two decompositions at slope zero}

Fix a finite-slope semistable $E$ at $(A,b)$ with $|b|\leq1$ and
$\nu_{A,b}(E)=0$.  Suppress $B=B_0+bH$ in numerical subscripts when
there is no ambiguity.  Write
\[
 V=\mathcal H^{-1}(E),\qquad U=\mathcal H^0(E),\qquad
 \bar U=U/U_{\rm tor}.
\]
The sheaf $V$ is torsion-free.  All the ordinary slope factors of
$V$ have slope at most $b-A$, and all those of $\bar U$ have slope
at least $b+A$.  To prove the first assertion, the highest slope
factor $V_1[1]$ is a subobject of $E$ in the first-tilt heart.
Finite-slope semistability excludes a denominator-zero subobject
and gives
\[
 w_B(V_1)\geq A^2r(V_1)/2.
\]
Classical Bogomolov--Gieseker gives
$2w_B(V_1)/r(V_1)\leq(\mu(V_1)-b)^2$; the first torsion pair gives
$\mu(V_1)\leq b$.  Together they give $\mu(V_1)\leq b-A$.
For the other assertion use the lowest slope factor of $\bar U$,
which is a quotient of $E$, and the same computation with
$\mu>b$.

Throughout this section, $O(e)$ denotes an absolute value bounded
by a fixed constant times $e=d-A|r|$.  Constants in this notation
may depend on the choices above, but not on $E$, its rank, or $b$.
Rank and character bounds for a complex do not mean bounds for the
length of its zero-dimensional cohomology.

For nonnegative rank, we retain a quotient with slopes within one
unit of $b+A$ and bound sections of the residual complex by $e$.
For nonpositive rank, we retain a shifted subsheaf with slopes
within one unit of $b-A$; its surface first cohomology will allow us to
bound all sections of $E$ by $e$.  This second bound will also be
applied to the shifted dual of a nonnegative-rank object, controlling
its degree-two cohomology.

Suppose first that $r(E)\geq0$.  On refined ordinary slope
filtrations the equality $e=d-Ar$ reads
\begin{equation}\label{eq:positive-excess}
 e=\sum_{\bar U}r_i(\mu_i-b-A)
   +d(U_{\rm tor})+\sum_Vr_i(b-\mu_i+A).
\end{equation}
Every term is nonnegative.  Let $Q$ be the slope tail quotient of
$\bar U$ whose slopes lie in $[b+A,b+A+1]$, and let $P$ be the kernel
of the epimorphism $E\to Q$ in the tilt heart:
\begin{equation}\label{eq:positive-residual-triangle}
 P\longrightarrow E\longrightarrow Q.
\end{equation}
Empty parts mean zero sheaves.  The negative cohomology of $P$ is
$V$, and its degree-zero cohomology is an extension of the slope
prefix with slopes $>b+A+1$ by $U_{\rm tor}$.  The corresponding
sheaf sequences, shifted where appropriate, are exact in the tilt
heart because their torsion-free terms stay in the indicated halves
of the torsion pair.

Equation~\eqref{eq:positive-excess} gives
\begin{equation}\label{eq:residual-character-bound}
 d(P),\ |r(P)|,\ r(\mathcal H^{-1}(P)),\
 r(\mathcal H^0(P))=O(e).
\end{equation}
For example, on the high prefix the positive number
$\mu_i-b-A>1$ controls both $r_i$ and $r_i(\mu_i-b)$; on $V$ the
number $b-\mu_i+A\geq2A$ does the same.  The torsion degree is at
most $e$.  Quotient comparison and the classical discriminant
inequality on the factors of $Q$ give
\begin{equation}\label{eq:positive-tail-moments}
 \frac{A^2r(Q)}2\leq w_B(Q)
 \leq\frac{A^2r(Q)}2+O(e),\qquad
 \sum_Qr_i\delta_i^B=O(e).
\end{equation}
Indeed, write $\mu_i-b=A+s_i$, where $0\leq s_i\leq1$ and
$\sum r_is_i\leq e$.  Then
$\sum r_i(\mu_i-b)^2=A^2r(Q)+O(e)$, and subtract $2w_B(Q)$.
Additivity with $w_B(E)=A^2r(E)/2$ now also gives $w_B(P)=O(e)$.

For $r(E)\leq0$ use instead
\begin{equation}\label{eq:negative-excess}
 e=\sum_Vr_i(b-A-\mu_i)
   +\sum_{\bar U}r_i(\mu_i-b+A)+d(U_{\rm tor}).
\end{equation}
Let $V_0$ be the slope prefix of $V$ with slopes in
$[b-A-1,b-A]$ and form the exact triangle in the tilt heart
\begin{equation}\label{eq:negative-residual-triangle}
 V_0[1]\longrightarrow E\longrightarrow P.
\end{equation}
The same estimates give \eqref{eq:residual-character-bound} and
$w_B(P)=O(e)$.  More explicitly, subobject comparison applied to
$V_0[1]$ gives
\[
 A^2r(V_0)/2\leq w_B(V_0)\leq A^2r(V_0)/2+O(e).
\]
Put $T=B-AH$.  If $\mu_i-b=-A-s_i$ with $0\leq s_i\leq1$,
then $\sum r_is_i\leq e$ and
\[
 \sum_{V_0}r_i\mu_T^2\leq e,\qquad
 w_T(V_0)=w_B(V_0)+A d_B(V_0)+A^2r(V_0)/2\geq-Ae.
\]
Since $\delta_i^T=\mu_T^2-2w_T(G_i)/r_i$, it follows that
\begin{equation}\label{eq:negative-tail-energy}
 \sum_{V_0}r_i(\mu_T^2+\delta_i^T)=O(e).
\end{equation}
After very general restriction to a surface in the fixed smooth
neighborhood and refinement into stable factors, this bound remains
valid by Lemma~\ref{lem:restriction-variance}.  For each such factor,
\[
 \mu_T^2+\delta^T=\|c_0/r-T\|^2+\xi.
\]
Consequently Proposition~\ref{prop:negative-surface-cohomology}
implies, for every integer $t_0<t\leq M$,
\begin{equation}\label{eq:negative-tail-h1}
 h^1\bigl(V_0|_S(t)\bigr)=O(e).
\end{equation}
Both decompositions are available when $r(E)=0$; we use the one
specified by the argument in which it occurs.

\subsection{Line comparisons and images of sections}

We next control the sections of $P$ in
\eqref{eq:positive-residual-triangle}, and those of $E$ in
\eqref{eq:negative-residual-triangle}.  Their restrictions can have
more sections than are bounded by $e$.  The argument will bound the
image of a space of sections under restriction.

We need a comparison valid for all tilt quotients of a line bundle.
At parameters $(a,b')$, let $\ell\in\ZZ$, $L=\OO_X(\ell H)$ and
$u=\ell-b'>0$.  A nonzero proper tilt quotient $I$ fits into
$0\to K\to L\to I\to0$ in the heart, with $K$ a sheaf and
\[
 0\longrightarrow\mathcal H^{-1}(I)\longrightarrow K
 \longrightarrow L\longrightarrow\mathcal H^0(I)\longrightarrow0.
\]
The image in $L$ is nonzero: otherwise $K=\mathcal H^{-1}(I)$
belongs to both halves of the slope torsion pair and is zero, so
the quotient is the whole line.  The image therefore has rank one.
It follows that $K$ is torsion-free, all its slopes lie in $(b',\ell]$,
and $r(I)=-r(\mathcal H^{-1}(I))$.  Classical Bogomolov--Gieseker
on the slope factors gives $w_{B_0+b'H}(K)\leq u d(K)/2$.
Since $r(L)=1$, $d(L)=u$ and
$2w(L)=u^2-\delta_{\rm line}$, subtraction gives
\begin{equation}\label{eq:line-quotient-numerator}
 n(I)\geq\frac{u d(I)-\delta_{\rm line}
                    +a^2r(\mathcal H^{-1}(I))}{2}.
\end{equation}
If $a^2\geq\delta_{\rm line}$, every finite quotient slope is at
least $u/2$ when $r(I)<0$, and at least
$u/2-h\delta_{\rm line}/2$ when $r(I)=0$.
For the latter assertion, $d(I)>0$ lies in $h^{-1}\ZZ$, so
$d(I)\geq1/h$.  The whole line has slope
\begin{equation}\label{eq:line-whole-slope}
 \nu_{a,b'}(L)=\frac u2-\frac{a^2+\delta_{\rm line}}{2u}.
\end{equation}

Choose the positive integer $m_0$ so large that
\[
 m_0-1>h\delta_{\rm line},\qquad
 (m_0-1)^2>A^2+\delta_{\rm line}.
\]
At $(A,b)$ all finite quotient slopes of $\OO_X(m_0H)$ are then
strictly positive.  A nonzero map from this line to $E$ would have
an image that is both such a quotient and a subobject of the
slope-zero semistable $E$.  A denominator-zero image is also
excluded by finite-slope semistability.  Thus, with $t_0=-m_0$,
\begin{equation}\label{eq:initial-section-vanishing}
 h^0(E(t_0))=0,\qquad
 h^0(P(t_0))=0\quad\hbox{in the positive-rank decomposition}.
\end{equation}

For the remaining estimates work at the side point
\[
 b_s=b-c,\qquad a_s=\sqrt{A^2-c^2},\qquad B_s=B_0+b_sH.
\]
Lemma~\ref{lem:tilt-semicircle} shows that $E$ is semistable there
with slope $c$.
Both triangles above remain short exact sequences in its tilt
heart: the negative sheaf slopes are at most $b-A<b-c$, and the
positive ones at least $b+A>b-c$.  The numerical quantities
$d(P),r(P),w(P)$ are still $O(e)$ after this bounded twist change.
For any integer $t_0\leq t\leq M$, the line $\OO_X(-tH)$ has
$u=-t-b_s>1$.  Equations~\eqref{eq:line-quotient-numerator} and
\eqref{eq:line-whole-slope} provide a single lower bound $-C_1$ for
all its finite tilt quotient slopes, uniformly in this finite range.

\begin{lemma}\label{lem:restriction-image}
Let $E\in\mathcal B_{H,B_0+bH}$ be finite-slope
$\nu_{A,b}$-semistable with $\nu_{A,b}(E)=0$ and $|b|\leq1$.
Put $B=B_0+bH$ and $e=d_B(E)-A|r(E)|$.  Let $P$ be the residual term in
\eqref{eq:positive-residual-triangle} when $r(E)\geq0$, or in
\eqref{eq:negative-residual-triangle} when $r(E)\leq0$.
Let
$W\subset\Hom(\OO_X(-tH),P)$ be any finite-dimensional subspace,
where $t_0<t\leq M$ is an integer.  There is a very general smooth
$S\in|H|$ in the fixed neighborhood, depending on this finite data,
such that the image of
\[
 W\longrightarrow\HH^0\bigl(S,P|_S(t)\bigr)
\]
has dimension at most $C_2e$.  The constant is independent of $W$
and its dimension.
\end{lemma}

\begin{proof}
Put $N=\dim W$ and take the tilt image $I$ and kernel $K'$ of the
evaluation map $\OO_X(-tH)^N\to P$ at the side point.
We first control the rank, first character, and $H$-degree of the
second character of $I$ by $e$, independently of $N$.  We then use
the evaluation kernel, whose rank is close to $N$, to cancel all
but $O(e)$ of the sections in the evaluation source.
Ordinary cohomology gives a sheaf $K'$ and an exact sequence
\begin{equation}\label{eq:section-evaluation-sheaves}
 0\longrightarrow V_I\longrightarrow K'
 \longrightarrow\OO_X(-tH)^N\longrightarrow U_I\longrightarrow0,
 \qquad V_I\subset V_P,
\end{equation}
where $V_I=\mathcal H^{-1}(I)$, $U_I=\mathcal H^0(I)$ and
$V_P=\mathcal H^{-1}(P)$.  Since $I\subset P$,
$0\leq d(I)\leq d(P)=O(e)$.  The sheaf $U_I(t)$ is globally
generated.  Its first Chern class has nonnegative $H$-degree,
including its divisorial torsion, and $H^2B_0=0$; hence
\[
 d(U_I)\geq u r(U_I),\qquad d(V_I)\leq0.
\]
Thus $r(U_I)=O(e)$.  The inclusion into $V_P$ gives
$r(V_I)=O(e)$ as well.

In the positive decomposition $I\subset P\subset E$, so
$n(I)\leq c d(I)$.  In the negative decomposition $P/I$ is a
quotient of $E$ and satisfies $n(P/I)\geq c d(P/I)$, also when its
denominator is zero.  Additivity and the numerical bounds on $P$
therefore give $n(I)\leq O(e)$ in both cases.

Every finite Harder--Narasimhan slope of $I$ is at least $-C_1$.
Indeed, its last factor receives a nonzero map from one of the
lines in the evaluation map.  The tilt image of this map is a
quotient of that line and a subobject of the last semistable
factor; the line comparison forces the claimed lower bound.
If the last factor has infinite slope there are no finite factors.
Denominator-zero factors have nonnegative numerator.  Writing
$d_i,n_i$ for the factors, we obtain
\[
 \sum_i|n_i|\leq n(I)+2C_1d(I)=O(e).
\]
Lemma~\ref{lem:tilt-support-bound}, applied at the fixed side
volume and summed over the filtration, now gives
\[
 \|c_{B_s}(I)\|+|w_{B_s}(I)|=O(e).
\]
The twist change to $I(t)$ is bounded.  Its rank is $O(e)$ by
\eqref{eq:section-evaluation-sheaves}, so
\begin{equation}\label{eq:evaluation-image-numerics}
 \|c_0(I(t))\|+|w_0(I(t))|=O(e).
\end{equation}

The sheaf $K'(t)$ is torsion-free and has rank $N+O(e)$.
Before twisting it belongs to the positive part of the torsion
pair, while \eqref{eq:section-evaluation-sheaves} bounds its maximum
slope by $-t$.  Its slopes after twisting therefore lie in
$(-u,0]$.  Its first and second characters are the negatives of
those in \eqref{eq:evaluation-image-numerics}.  For its refined
ordinary slope factors this proves
\begin{equation}\label{eq:evaluation-kernel-moments}
 \sum_i r_i\mu_i^2+\sum_i r_i\delta_i^0=O(e).
\end{equation}
In detail, the bounded nonpositive slope interval and the total
degree $O(e)$ bound the first sum; subtracting the total $2w_0$
bounds the second.

The next step is to obtain almost as many sections of this kernel
on a surface as the trivial evaluation source has.  Its small
first character and projected second character control the Riemann--Roch error;
the same numerical bounds will control the cohomology lost from
the negative term $V_I$.

Choose a very general $S$ in the fixed neighborhood and a very
general curve $C_m\in|mH|_S$ in its fixed neighborhood.
Impose also the finitely many open conditions that all the sheaf
sequences and filtrations used here restrict exactly, with
torsion-free restrictions where required.  Such choices are
possible by Lemma~\ref{lem:uniform-sections}.  Intersect the
refined restriction filtrations of $K'(t)$ with $V_I(t)$.
The successive intersections are subsheaves of the stable
restriction factors and have total rank $O(e)$.  By
Lemma~\ref{lem:restriction-variance}, the second slope moments
on both $S$ and $C_m$ are $O(e)$.  Applying
Lemma~\ref{lem:subsheaf-sections} to the intersections, after the
additional fixed twist $m$, yields
\begin{equation}\label{eq:evaluation-negative-sections}
 h^0\bigl(V_I|_S(t+m)\bigr)+
 h^0\bigl(V_I|_{C_m}(t+m)\bigr)=O(e).
\end{equation}
The reason this has no $N$ term is that each bound is a constant
times the rank of the intersection plus its rank times
$\mu_{i,+}^2$; its rank is at most $r_i$, and the second moments
in \eqref{eq:evaluation-kernel-moments} control the sum.

We also have $h^2(K'|_S(t+m))=O(e)$.  A stable factor of $K'(t)|_S$
with slope $\mu_i>\mu(K_S)-m$ has no homomorphism to $K_S(-m)$.
For all the other factors, \eqref{eq:apex-m-choice} gives
$\mu_i\leq-1$, so their total rank, as well as their second slope
moment, is $O(e)$.  Serre duality and
Lemma~\ref{lem:subsheaf-sections} applied to their dual hulls,
twisted by $K_S(-m)$, bound their $h^2$ by $O(e)$.
Surface Riemann--Roch and
\eqref{eq:evaluation-image-numerics} give
\[
 \chi\bigl(K'|_S(t+m)\bigr)
 =r(K')\chi(\OO_S(m))+O(e).
\]
Indeed the error consists of $w_0(K'(t))$ and the intersection of
$c_0(K'(t))$ with the fixed class $mH-K_S/2$; here
$K_S=(K_X+H)|_S$.  Consequently
\begin{equation}\label{eq:kernel-many-sections}
 h^0\bigl(K'|_S(t+m)\bigr)
 \geq r(K')h^0(\OO_S(m))-O(e).
\end{equation}

Restrict the triangle $K'\to\OO_X(-tH)^N\to I$ and twist by
$(t+m)H$.  The kernel on $H^0$ from $K'|_S(t+m)$ has dimension at
most
$h^{-1}(I|_S(t+m))=h^0(V_I|_S(t+m))$.
Equations~\eqref{eq:evaluation-negative-sections} and
\eqref{eq:kernel-many-sections}, with $r(K')=N+O(e)$, therefore
give
\[
 \begin{aligned}
 &\dim\im\!\left(
 H^0(K'|_S(t+m))\longrightarrow H^0(\OO_S(m))^N
 \right)\\
 &\hspace{25mm}\geq N h^0(\OO_S(m))-O(e).
 \end{aligned}
\]
The source of the next map has dimension $Nh^0(\OO_S(m))$.
Subtracting the displayed lower bound shows that the image of
\[
 H^0(\OO_S(m))^N\longrightarrow\HH^0(I|_S(t+m))
\]
has dimension $O(e)$.  This is the cancellation which removes
the unrestricted dimension $N$ from the estimate.

It remains to remove the auxiliary twist $m$.
Multiplication by the section defining
$C_m$ gives a map from $\HH^0(I|_S(t))$ to this latter group.
Its kernel has dimension at most
$h^{-1}(I|_{C_m}(t+m))$, which is again $O(e)$ by
\eqref{eq:evaluation-negative-sections}.  The image of
$H^0(\OO_S)^N$ in $\HH^0(I|_S(t))$ consequently has dimension
$O(e)$.  Composing with $I\to P$ proves the assertion for $W$.
All equalities involving $h^{-1}$ use derived restriction with
nonzerodivisors on the ordinary cohomology sheaves; no exactness
of a tilt-image operation under restriction was assumed.
\end{proof}

\begin{corollary}\label{cor:residual-sections}
Let $E\in\mathcal B_{H,B_0+bH}$ be finite-slope
$\nu_{A,b}$-semistable with $\nu_{A,b}(E)=0$ and $|b|\leq1$.
Put $B=B_0+bH$ and $e=d_B(E)-A|r(E)|$.  For $r(E)\geq0$, take $P$ from
\eqref{eq:positive-residual-triangle} and put
$V_P=\mathcal H^{-1}(P)$.  For $r(E)\leq0$, use the decomposition
\eqref{eq:negative-residual-triangle} and put $V=\mathcal H^{-1}(E)$.
There is a constant $C_3$, independent of $E,b$, such that
\begin{equation}\label{eq:residual-sections}
 \begin{aligned}
 h^0(P(j))&\leq C_3e &&\text{for the positive decomposition},\\
 h^0(E(j))&\leq C_3e &&\text{for the negative decomposition}.
 \end{aligned}
\end{equation}
Moreover
\begin{equation}\label{eq:negative-cohomology-vanishing}
 \begin{aligned}
 H^0(V_P(j_2))&=0&&\text{in the positive decomposition},\\
 H^0(V(j_2))&=0&&\text{in the negative decomposition}.
 \end{aligned}
\end{equation}
\end{corollary}

\begin{proof}
For the positive decomposition, apply
Lemma~\ref{lem:restriction-image} to all of $H^0(P(t))$ and the
triangle $P(t-1)\to P(t)\to P|_S(t)$.
The increase in $h^0$ is at most the dimension of the restriction image,
at most $C_2e$.  Telescope over the fixed set of integers
$t_0<t\leq j$, starting with
\eqref{eq:initial-section-vanishing}.

For the negative decomposition, apply the lemma to the image of
$H^0(E(t))$ in $H^0(P(t))$.  The kernel of
\[
 \HH^0(E|_S(t))\longrightarrow\HH^0(P|_S(t))
\]
has dimension at most $h^1(V_0|_S(t))=O(e)$ by
\eqref{eq:negative-residual-triangle} and
\eqref{eq:negative-tail-h1}.  Thus the restriction image of
$H^0(E(t))$ itself has dimension $O(e)$.  Telescope the restriction
triangle for $E$ from the same vanishing at $t_0$.
Finally every slope of the negative cohomology sheaves in
\eqref{eq:negative-cohomology-vanishing}, after twisting by $j_2$,
is at most $b-A+j_2<0$, by \eqref{eq:apex-A-choice}.
They have no sections.
\end{proof}

\subsection{Fixed sheaves in positive characteristic}

We have now obtained bounds involving only $e$ for the remainder
and for an object of nonpositive rank.  The positive tail $Q$ can
have arbitrarily large rank even when $e$ is small.  Its contribution
will instead be bounded by a cubic expression in its narrow range
of slopes.

Fix for the moment the individual object $E$.  All sheaves,
complexes, triangles, restriction maps and filtrations used for
this $E$ can be spread to a finitely generated integral subring
of $\CC$.  Localize so that the varieties have smooth geometrically
integral projective fibers, the finitely many coherent sheaves are
flat, and the displayed sequences and cohomology sheaves commute
with taking fibers.  Finite perfect presentations on the smooth
varieties justify the same assertion for the complexes.
Upper semicontinuity preserves the finitely many bounds
\eqref{eq:residual-sections} and vanishings
\eqref{eq:negative-cohomology-vanishing} after further localization.
We do not require tilt semistability on the reductions.

We will require ordinary slope semistability for finitely many
fixed sheaves on $X$, on chosen surfaces, and on chosen curves.
For completeness, the needed spreading assertion and its
Frobenius consequence are recorded next.

\begin{lemma}\label{lem:fixed-sheaf-spreading}
Let $G$ be a slope-semistable torsion-free sheaf on a smooth
projective complex variety $Y$ with a fixed very ample class.
There is a model and a localization on which the geometric fibers
of $G$ are torsion-free and slope-semistable.  Finitely many such
requirements may be imposed simultaneously.
\end{lemma}

\begin{proof}
Embed $G$ in $\OO_Y(u)^k$ by choosing generators for a twist of
its dual and using the injection $G\to G^{**}$.  Spread this
injection and make its cokernel flat, so it remains an injection
on the fibers.  We describe a finite-type set containing every
destabilizing saturated subsheaf on every geometric fiber.

For a saturated subsheaf $J$ of rank $s$, $0<s<r(G)$, its generic
Pl\"ucker coordinates extend across codimension two to sections
of $\OO_Y(su)\otimes(\det J)^{-1}$.  If $J$ destabilizes, the
degree of the zero divisor of a nonzero coordinate is bounded
above by a number depending only on $u,s$ and the degree of $G$.
An effective divisor of degree at most $d_0$ is contained, with
its multiplicities, in the zero divisor of a section of
$\OO_Y(v)$ with $v\leq\max(1,d_0)$.  To see this, choose a finite
linear projection to projective space of dimension $\dim Y$.
Each prime component has hypersurface image of degree at most
its own degree.  Pull back an equation for each such image and
multiply with the required multiplicities.  The projection is
finite because its pullback of $\OO(1)$ is ample.

Divide the Pl\"ucker tuple by one nonzero coordinate, then multiply
by this section to clear its pole divisor.  Normality extends the
resulting tuple to global sections of $\OO_Y(v)$.  The saturated
subsheaf $J$ is the kernel, inside $G$, of wedging with this tuple,
viewed as a map to a fixed sum of $\OO_Y(u+v)$: the kernels agree
at the generic point and are saturated.  The finitely many possible
values of $s,v$ give finite-type spaces of tuples after shrinking
the base so that their section spaces commute with base change.
On each parameter space, flatten the cokernel of the universal
wedge map.  The kernel then commutes with taking fibers, and its
rank and degree are constructible, as read from its Hilbert
polynomial.  Hence the condition that some tuple gives a kernel
of smaller positive rank and larger slope than $G$ is constructible
on the base.  Allowing tuples not satisfying the Pl\"ucker equations
does not cause a problem: a kernel with those numerical properties
is itself a destabilizing subsheaf.

This constructible subset excludes the generic point.  Otherwise
a destabilizer over its algebraic closure would extend to one
over $\CC$.  A constructible subset of an integral noetherian
scheme whose closure is the whole scheme contains a nonempty
open set and hence the generic point.  Its closure is therefore
proper; removing that closure proves the assertion.
\end{proof}

\begin{lemma}\label{lem:frobenius-slope-gap}
On a smooth projective geometric fiber in characteristic $p$, let
$G$ be a slope-semistable torsion-free sheaf of rank $r$.
For a constant $c_Y$ fixed by the chosen model of $Y$ and its
polarization,
\begin{equation}\label{eq:frobenius-slope-gap}
 \mu_{\max}(F^*G)\leq p\mu(G)+c_Y(r-1).
\end{equation}
\end{lemma}

\begin{proof}
We use the Cartier-descent argument of
\cite[Theorem~4.1 and Proposition~4.2]{Langer2022}, keeping the
chosen cotangent twist in the slope normalization.
Generation of a fixed positive twist of the tangent bundle gives
an embedding $\Omega_Y^1\hookrightarrow\OO_Y(c_Y)^k$ with locally
free cokernel.  It can be chosen on the model and retained on its
fibers.  Frobenius is finite flat because the fiber is smooth over
a perfect field.

Suppose two consecutive Harder--Narasimhan slopes of $F^*G$ differ
by more than $c_Y$.  For the truncation $J$ above that gap, the
second fundamental map for the canonical connection is
\[
 J\longrightarrow(F^*G/J)\otimes\Omega_Y^1.
\]
It is $\OO_Y$-linear.  Its source has minimum slope larger than
the maximum slope of its target, the latter bounded using the
displayed embedding, so the map is zero.  The connection thus
preserves $J$.  The canonical connection has zero $p$-curvature,
as does its restriction to $J$.  Cartier descent therefore
descends this inclusion to a subsheaf of the Frobenius twist of
$G$.  Its slope is $\mu(J)/p>\mu(G)$, contradicting semistability.
Here the Frobenius twist of the algebraically closed base field
is an automorphism and does not change the slope assertion.
Thus no adjacent gap is greater than $c_Y$.  There are at most
$r$ factors, so the maximum slope is at most the mean plus
$c_Y(r-1)$, as claimed.
\end{proof}

The localizations above concern finitely many fixed sheaves, after
$E$ has been chosen.  The resulting base still has geometric closed
fibers in unbounded prime characteristics: its constructible image
in $\Spec\ZZ$ contains the generic point and hence a nonempty open
set.  Very-generality was used only to choose the original complex
divisors; it is not asserted for their reductions.

\subsection{The tail and the Frobenius limit}

Suppose $r(E)\geq0$ and use
\eqref{eq:positive-residual-triangle}.  For every stable ordinary
slope factor $G$ of $Q$, choose a very general smooth
$S\in|H|$ and a very general $C\in|H|_S$, imposing all required
exactness conditions.  Refine $G|_S$ into stable factors $G_i$
and their restrictions to $C$ into stable factors $G_{ij}$.
These curves, which serve only in the present algebraic argument,
need not belong to the neighborhoods used for the analytic
estimates.  Lemma~\ref{lem:restriction-variance} twice gives,
with $\mu=\mu(G)$,
\begin{equation}\label{eq:twice-restricted-variance}
 \sum_{i,j}r_{ij}(\mu_{ij}-\mu)^2
 \leq3r(G)\delta_G^B.
\end{equation}
Indeed the first variance is at most $r(G)\delta_G^B$.
The sum of $r_i\delta_i^B$ on $S$ equals $r(G)\delta_G^B$ plus
that variance, so is at most $2r(G)\delta_G^B$; apply the
restriction estimate once more and add the two variances.
The weighted means are preserved at each step.  Spread this
finite data and apply Lemma~\ref{lem:fixed-sheaf-spreading} to
$G,G_i,G_{ij}$.

Since $H^2D=0$, the normalized slope shift of $L_p$ on $X,S,C$ is
$-pb+O(1)$, with a bounded rounding error.  Frobenius commutes with
these restrictions, and its flatness preserves their exact
filtrations.  In the following estimates an $O_G$ constant may
depend on all the fixed data chosen for $G$, but not on $p$ or on
the fiber.  In particular the ranks in the Frobenius slope error
are now fixed, so they may enter these constants.  They must not
enter the coefficients which survive division by $p^3$.
Lemma~\ref{lem:frobenius-slope-gap} gives a vanishing
twist
\[
 k_X=-p(\mu-b)+O_G(1),\qquad
 H^0(F^*G\otimes L_p(k_X))=0,
\]
where the bounded integer adjustment is chosen to make the maximum
slope negative.  Telescope from $k_X$ to zero using the surface
restriction sequence.  At each integer $k$ in this interval,
bound the surface sections by those of its factors $G_i$.
For each such factor telescope on $S$ from its own vanishing
twist $-p(\mu_i-b)+O_G(1)$ up to $k$, if $k$ exceeds that twist.

On $C$, the $H$-degree is $h$.  Taking successive jets at a point
until the maximum ordinary degree slope is negative, and using
\eqref{eq:frobenius-slope-gap}, gives for every factor
\begin{equation}\label{eq:curve-ramp-bound}
 h^0(F^*G_{ij}\otimes L_p(l))
 \leq h r_{ij}\bigl(p(\mu_{ij}-b)+l\bigr)_++O_G(1).
\end{equation}
The error includes the integer rounding and at most a bounded
number of point jets, each of dimension $r_{ij}$; all ranks here
are fixed before $p$ varies.  Each surface telescope has only
$O_G(p)$ terms in the range under consideration.  Summing the
ramp in \eqref{eq:curve-ramp-bound}, and enlarging the sum down
to all its positive terms when necessary, gives
\begin{align}
 h^0(F^*G|_S\otimes L_p(k))
 &\leq\frac{hp^2}{2}\sum_{i,j}r_{ij}
        (\mu_{ij}-b+k/p)_+^2+O_G(p)\notag\\
 &\leq\frac{hp^2}{2}\bigl[
 r(G)(\mu-b+k/p)_+^2+3r(G)\delta_G^B\bigr]+O_G(p).
 \label{eq:surface-square-bound}
\end{align}
For the second inequality use
$x_+^2\leq y_+^2+2y_+(x-y)+(x-y)^2$ and the vanishing of the
weighted first displacement from the mean, followed by
\eqref{eq:twice-restricted-variance}.  The ramp sum differs from
the corresponding half-square by $O_G(p)$, uniformly for the
indices used above.

The outer telescope has length
$p(\mu-b)+O_G(1)\leq(A+1)p+O_G(1)$.  Summing
\eqref{eq:surface-square-bound} and using the Riemann sum for
the square on $[-(\mu-b),0]$ gives the following bound.
Here and below, limits are taken along geometric fibers of the chosen
model, after all the specified localizations, with characteristics
$p$ tending to infinity:
\begin{equation}\label{eq:tail-factor-frobenius}
 \limsup_{p\to\infty}p^{-3}h^0(F^*G\otimes L_p)
 \leq\frac{h r(G)}6(\mu-b)^3
       +C_Ah r(G)\delta_G^B.
\end{equation}
The summed $O_G(p)$ errors are $O_G(p^2)$ and vanish in this
limit.  One may take a fixed multiple of $A+1$ for $C_A$;
it does not depend on the restriction sheaves or their ranks.
Summing over the finite filtration of $Q$, using
\eqref{eq:positive-tail-moments}, gives
\begin{equation}\label{eq:tail-frobenius}
 \limsup_{p\to\infty}p^{-3}h^0(F^*Q\otimes L_p)
 \leq\frac{hA^2}{6}d_B(Q)+O(e).
\end{equation}
To verify its leading term, write $\mu-b=A+s$, $0\leq s\leq1$.
Then $(A+s)^3-A^2(A+s)\leq C_As$, and
$\sum r_is_i\leq e$ by \eqref{eq:positive-excess}.

It remains to bound the residual contribution and the positive
even cohomology group in degree two.  Let $P_1$ be the positive
residual $P$ or a nonpositive-rank object treated by the negative
decomposition.  It has ordinary cohomology in degrees $[-1,0]$,
and \eqref{eq:negative-cohomology-vanishing} supplies the vanishing
required in Corollary~\ref{cor:frobenius-section-transfer}.
That corollary and \eqref{eq:residual-sections} give, for either
rounding line $L$,
\begin{equation}\label{eq:residual-frobenius}
 h^0(F^*P_1\otimes L)
 \leq N_{p,L}h^0(P_1(j))\leq C_4p^3e.
\end{equation}
The characteristic-zero section bounds pass to these fibers by
the finite semicontinuity conditions already imposed.

Apply Lemma~\ref{lem:tilt-duality} to $E$:
\begin{equation}\label{eq:apex-duality-triangle}
 \widetilde E\longrightarrow R\mathcal Hom(E,\OO_X)[1]
 \longrightarrow T_0[-1],
\end{equation}
where $T_0$ is zero-dimensional and $\widetilde E$ is finite-slope
semistable of slope zero at $(A,-B)$.  Its rank is $-r(E)$,
its denominator is $d(E)$, its excess is $e$, and
$z_{-B}(\widetilde E)=z_B(E)+\operatorname{length}(T_0)/h$.
Include this triangle and the negative decomposition for
$\widetilde E$ in the fixed spreading data.  Serre duality,
flatness of $F$, and \eqref{eq:apex-duality-triangle} give
\begin{equation}\label{eq:degree-two-frobenius}
 h^2(F^*E\otimes L_p)
 =h^0(F^*\widetilde E\otimes L'_p)\leq C_4p^3e.
\end{equation}
For clarity, Serre duality first identifies the dual of the
degree-two group with degree zero of
$F^*R\mathcal Hom(E,\OO_X)[1]\otimes L'_p$.
The term $F^*T_0[-1]\otimes L'_p$ has zero hypercohomology in
degrees $-1$ and $0$, so it does not alter that group.

The complex $F^*E\otimes L_p$ has hypercohomology in degrees
$[-1,3]$.  The only even degrees in this interval are zero and
two, whence
\[
 \chi(F^*E\otimes L_p)
 \leq h^0(F^*E\otimes L_p)+h^2(F^*E\otimes L_p).
\]
The degree-zero group is bounded above by the sum of the
corresponding groups for $P$ and $Q$ in
\eqref{eq:positive-residual-triangle}.  Equations
\eqref{eq:tail-frobenius}, \eqref{eq:residual-frobenius} and
\eqref{eq:degree-two-frobenius} therefore imply
\begin{equation}\label{eq:euler-frobenius-bound}
 \limsup_{p\to\infty}p^{-3}\chi(F^*E\otimes L_p)
 \leq\frac{hA^2}{6}d_B(Q)+C_5e.
\end{equation}

Riemann--Roch identifies the limit on the left:
\begin{equation}\label{eq:rr-frobenius-limit}
 \lim_{p\to\infty}p^{-3}\chi(F^*E\otimes L_p)=h z_B(E).
\end{equation}
Here is a way to compare varying characteristics without
identifying their cohomology rings.  On a smooth fiber,
$\ch_i(F^*E)=p^i\ch_i(E)$ by the splitting principle.  The same
formula holds for the Adams class $\psi^p[E]$ in vector-bundle
$K$-theory.  Riemann--Roch thus identifies the Euler characteristic
in \eqref{eq:rr-frobenius-limit} with that of
$\psi^p[E]\otimes L_p$.  This class is defined by exterior-power
operations on a fixed perfect presentation and commutes with
specialization.  Its Euler characteristic is constant from the
model to the complex fiber.  On that fiber, write
$c_1(L_p)=-pB+R_p$, where $R_p$ remains in a fixed bounded set of
divisor classes.  The degree-three leading term of
\[
 \int_X\left(\sum_{i=0}^3p^i\ch_i(E)\right)
                  e^{-pB+R_p}\td(X)
\]
is $p^3\ch_3^B(E)$; all other terms are $O_E(p^2)$.
This proves \eqref{eq:rr-frobenius-limit}, without a condition on
$K_X$.

Combining \eqref{eq:euler-frobenius-bound} and
\eqref{eq:rr-frobenius-limit}, and using $d_B(Q)\leq d_B(E)$,
proves \eqref{eq:apex-estimate} for $r(E)\geq0$.
The surviving constants came only from the fixed regularity,
surface estimates and the interval $[A,A+1]$; every constant
depending on the individual spread sheaves occurred in an error
of order at most $p^2$ and disappeared after division by $p^3$.
This establishes the asserted order of quantifiers.

If $r(E)<0$, apply the just-proved case to $\widetilde E$ in
\eqref{eq:apex-duality-triangle}, which has positive rank.
Its $d$ and $e$ agree with those of $E$, and its third character
is at least $z_B(E)$.  The estimate for $\widetilde E$ therefore
implies the estimate for $E$.  Our constants were chosen for both
signs of $B_0$, so this use of the mirror parameters costs nothing.
Finally undo the integral tensor normalization of $b$.
This completes the proof of Proposition~\ref{prop:apex}.
